\documentclass[12pt]{article}
\usepackage[utf8]{inputenc}
\usepackage[T1]{fontenc}
\usepackage[margin=1in]{geometry}
\usepackage{setspace}
\usepackage{xcolor}
\usepackage[colorlinks=true, urlcolor=blue!60!black, linkcolor=blue!60!black]{hyperref}
\setlength{\parindent}{0em}
\setlength{\parskip}{0.45em}
\renewcommand{\baselinestretch}{1.0}
\begin{document}
\begin{center}
{\Large \textit{United States v. E.C. Knight Co.}, 156 U.S. 1 (1895)\par}
\vspace{0.3em}
{\normalsize Case Brief --- One-Page Edition\par}
\end{center}

\textbf{Memory Jogger.} The Sugar Trust case and the manufacture/commerce line (\href{https://supreme.justia.com/cases/federal/us/156/1/}{156 U.S. 1}).

\textbf{Facts.} In 1892 the American Sugar Refining Company bought four Philadelphia refineries, including the E.C. Knight Company, giving the ``Sugar Trust'' about 98\% of American sugar refining. The United States sued under the Sherman Act to cancel the purchases.

\textbf{Procedural History.} The Circuit Court dismissed the bill as concerning only manufacturing in Pennsylvania; the Third Circuit affirmed.

\textbf{Votes.} Affirmed, 8--1. Fuller wrote for the Court; Harlan dissented alone.

\textbf{Holding.} The commerce power, and so the Sherman Act, does not reach a monopoly of manufacture; manufacturing affects interstate commerce only incidentally and indirectly and is left to the states' police power.

\textbf{Reasoning.} Manufacture is transformation; commerce is buying, selling, and transportation, and the two spheres are exclusive. Otherwise little would be left for the states. Harlan answered that a trust controlling 98\% of a necessity sold nationwide directly burdens interstate commerce and that no state could stop it. \textit{Swift}, \textit{Jones \& Laughlin}, and \textit{Wickard} later repudiated the line.

\textbf{Examples.} \emph{Same-Side:} purely local, noneconomic possession (\textit{Lopez}). \emph{Opposite-Side:} a 95\% chip-plant merger today. \emph{Fence-Sitter:} homemade guns within a market scheme.

\textbf{Critique.} A study in formalism divorced from economic reality that forced later Courts to seek workable limits.

\textbf{Key Quotation.} ``Commerce succeeds to manufacture, and is not a part of it. \ldots That which belongs to commerce is within the jurisdiction of the United States, but that which does not belong to commerce is within the jurisdiction of the police power of the state.'' --- Chief Justice Fuller, Opinion of the Court, 156 U.S. at 12.

\end{document}
